\section{Conclusion}

As AI agents take on greater operational responsibility in defending critical
infrastructure, frameworks for understanding their failure behaviour become
essential, not as an afterthought, but as a prerequisite for responsible
deployment. \ARES{} provides an integrated treatment of this problem
for the cybersecurity domain: a failure taxonomy that classifies how agents
fail across origin, manifestation, and impact dimensions; a quantitative
reliability model based on skill--knowledge--capability mismatch; an open
evaluation testbed supporting controlled experimentation; and a labelled
failure dataset enabling reproducible benchmarking.

Our empirical results establish three findings with direct implications for
SOC operations. First, knowledge gaps are the dominant driver of
hallucination: agents operating without relevant threat context fabricate
information at alarming rates, making knowledge completeness verification
the highest-priority deployment check. Second, multi-agent pipelines amplify
rather than mitigate component failures, challenging the intuition that
task decomposition improves outcomes. Third, the Epistemic Gap (capturing
what an agent does not know it does not know) decomposes
false-negative risk into a knowledge term and a calibration term that
are statistically equivalent in raw predictive power but independently
controllable, making confidence calibration as actionable a lever as
knowledge coverage.

These findings translate into actionable deployment guidance: pre-execution
reliability screening via ARS, runtime confidence monitoring via the
Epistemic Gap, evidence-based knowledge refresh scheduling via the temporal
drift model, and conservative multi-agent pipeline design informed by
cascade reliability analysis.

\ARES{}, \ARESBench{}, and \SAFK{} are publicly released to support
reproducible research in AI security agent
dependability.\footnote{\url{https://github.com/IntelEyes/ARES-Framework}}
